% 来源: 19c9fbbe-eb32-8796-8000-0000e5e98b5d_4.2_数据拷贝.pdf
% 提取时间: 2026-02-28T22:51:37+08:00

4.2 数据拷贝
除了传统操作系统结构的限制，应用层与操作系统缓冲区之间的数据拷贝是可以避免的。消除这些冗余
拷贝有助于软件更好地发挥硬件性能潜力。本章将探讨如何通过技术手段降低数据操作成本，同时保持
系统模块化设计并避免大规模改动。
具体内容包括：
 1. 分析冗余数据拷贝的成因与机制。
 2. 通过重构操作系统和网络代码，在端节点本地避免数据拷贝。
 3. 利用远程直接内存访问（RDMA）技术，减少大传输时的拷贝与控制开销。

 4. 在Web服务器文件系统中消除文件缓存与应用间的冗余拷贝。
 5. 探讨校验和计算、加密等全量数据操作，提出分层集成处理方法。

 6. 超越数据拷贝问题，分析缓存使用不当对性能的影响及优化策略。


1 为何会出现数据拷贝
为展示处理服务器上的GET请求时涉及的冗余拷贝，下图描述在最坏情况下从磁盘读取文件数据到通过
网络适配器发送相应数据段所涉及的数据传输过程。
图中的主要硬件组件包括CPU、内存总线、I/O总线、磁盘和网络适配器。主要的软件组件包括Web服务
器应用程序和内核。涉及的两个主要内核子系统分别是文件系统和网络系统。为了简化模型，图中仅展
示了一个CPU（许多服务器是多处理器，尤其是多核机器），并且只关注静态内容的请求（许多请求是
动态内容，由通用网关接口（CGI）进程提供，通过某种进程间通信机制，如UNIX管道，将这些内容传
递给服务器）。




                                                1
直观地说，故事很简单。文件通过read()系统调用从磁盘读取到应用程序缓冲区。然后，HTTP响应与应
用程序缓冲区的内容一起通过write()系统调用经TCP连接发送到客户端。网络子系统中的TCP代码将响
应数据分割成适合传输的数据段，并在每个数据段中添加TCP校验和后将其传输到网络适配器。
真实情况会稍微复杂。首先，文件通常被读取到内核内存中的一个区域，称为文件缓存（图中的
copy1）。这样做的好处是，后续对热门文件的请求可以直接从主内存中提供服务，而无需缓慢的磁盘
I/O。然后，文件从文件缓存拷贝到Web服务器应用程序缓冲区（图中的copy2）。由于应用程序缓冲区
和文件缓存缓冲区位于主内存的不同部分，这种拷贝只能通过CPU从第一个内存位置读取数据并写入第
二个内存位置来完成。
随后，Web服务器执行write()到相应的套接字。由于应用程序可以在write()之后随时重用其缓冲区（甚
至释放它），网络子系统中的内核不能简单地从应用程序缓冲区传输数据。特别是TCP软件可能需要在
不可预测的时间后重新传输文件的一部分，到那时应用程序可能希望将其缓冲区用于其他目的。
因此，UNIX（以及许多其他操作系统）提供了所谓的拷贝语义。write()调用中指定的应用程序缓冲区被
拷贝到一个套接字缓冲区（内核中的另一个缓冲区，位于与文件缓存或应用程序缓冲区不同的内存地
址）。这在图中称为copy3。最后，每个数据段在添加IP和链路头后被发送到网络（通过拷贝数据从套接
字缓冲区到网络适配器内的内存），这称为copy4。
在这四个拷贝操作和校验和计算中，每个操作都会消耗资源。所有四个拷贝操作以及校验和计算都会消
耗内存总线的带宽。内存位置之间的拷贝（copy2和3）实际上更糟糕，因为它们需要对总线进行一次读
和一次写操作来传输内存中的每个bit。TCP校验和需要对每个bit进行一次读操作和一次写操作来附加最
终的校验和。最后，如果CPU执行繁重的拷贝工作（所谓的编程I/O），copy1和4的开销可以与copy2和
3一样高；然而，如果设备本身执行拷贝（所谓的DMA），仅需要一次总线读或写操作。
这些拷贝操作还会消耗I/O总线带宽，最终消耗内存带宽本身。假设内存每 x 纳秒提供一个大小为 W
位的字，则内存的理论吞吐量极限为 W /x 位每纳秒。例如，即使假设使用DMA，这些拷贝操作确保服
                                                     2
务器的内存总线被每个文件发送的字使用7次。因此，Web服务器的吞吐量不能超过 T /7 ，其中 T 是
内存和内存总线中较慢者的速度。其次，额外的拷贝操作消耗了更多的内存。同一个文件可以存储在文
件缓存、应用程序缓冲区和套接字缓冲区中。
综上所述，冗余拷贝以两种基本且相互独立的方式对性能造成损害。首先，由于占用了超出实际需求的
总线和内存带宽，Web服务器的运行效率受到制约，其性能表现低于总线的理论传输能力，即便在处理
驻留在内存中的文档时亦是如此。其次，由于内存资源的过度消耗，Web服务器被迫从磁盘读取数量显
著高于正常比例的文件，而非充分利用文件缓存机制，从而进一步加剧了性能瓶颈。
理想情况下，所有这些繁琐的总线传输操作都应被消除。显然，如果数据已在缓存中，则无需进行
copy1，因此可以忽略它。如果数据不在缓存中，服务器将以磁盘速度运行，这本身就太慢了（尽管使用
快速的非易失性存储器会改变这一状况）。copy3似乎没有必要——为什么不能直接将数据从文件缓存的
内存位置发送到网络呢？同样，copy2也似乎没有必要。而拷贝4则是不可避免的。
2 通过本地重构减少拷贝
本节首先集中讨论Copy 3，即当网络消息被发送（接收）时，从应用程序到内核缓冲区（或反之）的基
本拷贝。这是网络中的一个基本问题，独立于文件系统问题。假设协议是固定的，但本地实现（至少是
内核）可以重构。当然，目标是进行最小的重构，以便继续利用现有内核和应用程序软件中的大量投
资。我们首先描述基于利用适配器内存的技术。然后，描述拷贝避免的核心思想（通过共享物理页面的
重新映射）及其陷阱。其次，展示如何通过基于I/O流的预计算和缓存来优化页面重新映射；然而，这项
技术涉及更改应用程序编程接口（API）。最后，第4节描述了另一种技术，该技术使用虚拟内存
（VM），但不更改API。
2.1 利用适配器内存
在内存映射架构中，内存可位于总线上的任意位置。内存映射机制表明，CPU通过读写设备物理内存空
间的某一部分，与所有外围设备（如适配器和磁盘）进行通信。因此，尽管内核内存通常驻留在主内存
子系统中，理论上并无限制阻止我们将部分内核内存部署在适配器上，而适配器通常具备一定的内存资
源。
通过充分利用适配器的现有内存，并利用这种灵活的内核内存布局策略，可将部分内核内存直接放置于
适配器上。这一设计的优势在于，当数据从应用程序复制到内核内存后，实际上已存在于适配器中，从
而避免了为完成网络传输而进行的额外数据拷贝操作。最终，这种优化不仅减少了数据传输的开销，还
显著提升了系统效率。如下图所示：




                                                  3
                消除内核到适配器之间的拷贝需求
注意图中忽略了任何从磁盘到内存的传输。本质上，无用的Copy 3现在与必要的Copy 4结合，形成一次
拷贝。
校验和的机制主要基于费用分摊的理念。其核心思想是，在数据从应用程序缓冲区传输到内核内存（由
适配器驻留）的过程中，处理器通过编程I/O（PIO，即由处理器主导的输入/输出操作）读取数据包的每
个字。由于这种总线读取操作具有较高的开销，CPU可以在数据拷贝过程中顺带计算校验和。具体而
言，CPU通过维护一个累加寄存器来存储传输字的累加和，从而实现这一目标。
然而，实际的数据传输任务并非由CPU直接完成，而是由适配器通过直接内存访问（DMA）机制来执
行。因此，当CPU不再参与数据拷贝过程时，校验和的计算责任应转移到适配器上，并在DMA传输过程
中对数据字进行校验和计算。
尽管这种设计在理论上具有一定的吸引力，但它存在三个显著的缺陷。首先，这种方法要求网络适配器
具备大量内存以支持多个高吞吐量的TCP连接（这些连接通常需要较大的窗口尺寸）。这将导致适配器
的成本显著增加，可能超出用户的预期预算。其次，在某些实现方式中，校验和由CPU在接收数据包并
将其拷贝到应用程序缓冲区时计算。然而，这种操作可能导致数据损坏的风险。虽然校验和不匹配的情
况最终会被检测到，但在此期间可能会引发尴尬的局面，例如应用程序可能已经读取了错误的数据。最
后，这种方式在性能和可靠性之间的权衡上存在一定局限性，可能无法完全满足高可靠性和高效率的系
统需求。
2.2 使用写时拷贝（copy-on-write）
我们来探讨一种新的思路，通过虚拟内存（VM）重映射技术，在大多数场景下避免应用程序和内核之间
的数据拷贝操作。这种技术的核心思想是利用写时拷贝（Copy-On-Write, COW）机制。回想一下，为
                                                    4
什么我们需要单独的数据拷贝？原因之一是应用程序可能会修改缓冲区的内容，从而破坏TCP协议的语
义；另一个原因是应用程序和内核运行在不同的虚拟地址空间中。
一些操作系统（比如Mach）提供了一种非常巧妙的机制——写时拷贝（COW）。借助这种机制，进程
可以在内存中以极低的成本复制虚拟页（VP）。具体来说，COW并不会真的立即复制整个页面的数据，
而是让新创建的副本和原始页面共享同一个物理页（P）。这只需要更新少量的描述符（也就是几个内存
字），而不是直接拷贝整个数据包（例如1500字节的数据）。只有当原始数据的所有者尝试修改数据
时，操作系统才会介入，自动生成两个独立的物理副本（P 和 P'）。此时，原始所有者继续使用 P 并可
以对其进行修改，而副本则指向旧的物理页 P'。虽然在这种情况下会触发物理页的完整拷贝，但这种情
况通常很少发生（或者只有少数页面会被修改），因此整体开销仍然较低。
基于这一机制，应用程序可以为内核创建一个写时拷贝的副本。如果应用程序在极少数情况下修改了缓
冲区，内核会执行一次昂贵的物理拷贝操作。然而，由于这种修改并不常见，系统整体性能依然得到了
优化。实际上，这正是惰性求值（Lazy Evaluation）的一种体现，旨在最小化常规情况下的开销。最
后，值得一提的是，校验和的计算可以通过CPU在数据传输到适配器内存或从适配器内存传回时完成，
也可以通过适配器上的CRC硬件来加速。
不过需要注意的是，许多常见的操作系统（如UNIX和Windows）并未直接提供写时拷贝的支持。但通过
理解COW背后的核心原理，即虚拟内存的灵活运用，我们仍然能够实现类似的效果，并从中受益！
实现copy-on-write




                    使用copy-on-write
现在大多数电脑都用虚拟内存（VM）技术。简单来说，程序员写程序时，可以假装有一个无限大的内存
空间。实际上，这个“无限内存”是虚拟的，编译器会把程序里的变量分配到一个虚拟地址表里。比如，
某个变量在虚拟地址表里的位置是1010，但这个地址并不直接对应真实的物理内存。虚拟地址会通过一
个叫页表的东西，映射到实际的物理内存上，而这些物理内存可能在硬盘或者主内存里。
                                                   5
虚拟地址分成两部分：高位的部分表示“页号”，低位的部分表示“页内偏移”。举个例子，假设虚拟地址
有32位，其中高20位是页号，低12位是页内偏移。物理内存也被分成一块一块的，每块叫一个物理页，
比如每 212 字节算一页。当程序访问一个虚拟地址时，系统会根据页号查页表，找到对应的物理页。如
果需要的页不在内存里（可能还在硬盘上），硬件就会发出一个信号，让操作系统把它从硬盘加载到内
存。
为了加快查找页表的速度，避免每次都去内存里找页表，还有一个叫TLB（Translation Lookaside
Buffer，转译后备缓冲器）的小缓存，专门用来快速找到虚拟地址对应的物理页。
说到这里，就能解释COW（写时复制）是怎么工作的了。假设有一个虚拟页X，它指向物理内存中的页
P。如果系统想把这个页的内容复制到一个新的虚拟页Y（比如另一个进程需要用到），最笨的方法是重
新分配一个物理页P'，然后把P的内容拷贝过去，再让Y指向P'。但其实有更聪明的办法：直接修改页
表，让新的虚拟页Y也指向原来的物理页P。因为现代操作系统通常用很大的页，改页表比拷贝数据快得
多。
不过，这里有个细节：为了让COW生效，系统会在原始虚拟页X的页表条目里设置一个特殊的标记，叫
COW保护位。如果程序想修改X的内容，硬件会发现这个标记，然后发出信号给操作系统。这时，操作
系统才会真正分配一个新的物理页P'，把P的内容拷贝过去，然后让X指向P'，同时清除COW标记。而虚
拟页Y仍然指向原来的物理页P。
2.3 Fbufs：优化页面重映射




上图给出了一个具体的例子来说明页面重映射的过程。假设操作系统想快速把进程1（比如某个应用程
序）中的数据复制到进程2（比如内核）的某个虚拟页（比如VP 8）。乍一看，似乎只需要修改进程2中
VP 8对应的页表条目，让它指向进程1中VP 10已经指向的数据就行了。然而，这个简单的说法忽略了很
多额外的开销和细节。

                                                      6
1.   多级页表：现在的系统大多用多级页表来做内存映射，因为如果把虚拟地址（比如20位的虚拟页
     号）直接映射到一个页表里，会占用太多内存。所以，实际操作时需要分几层页表来完成映射。为
     了方便在不同机器上使用，还有通用的和特定机器的页表设计。这样一来，修改页表就不只是简单
     的一次操作，而是要涉及多个步骤。
  2. 获取锁并修改页表条目：页表是大家共享的东西，为了避免冲突，必须加个“锁”来保护它。修改之
     前要先拿到这个锁，改完后再释放它。
  3. 刷新TLB缓存：为了加快地址转换速度，常用的页表映射会被存在一个叫TLB的缓存里。当你为某个
     虚拟地址（比如VP 8）写入新的位置时，必须找到TLB里对应的旧记录，并把它清除或更新，否则
     可能会出问题。
  4. 分配目标页的位置：虽然我们假设目标页的位置是VM 8，但在真正拷贝数据之前，还得先找一找目
     标进程里有没有空闲的页表条目。这一步也需要一些计算。
  5. 锁定页面防止被换出：物理内存有限，某些页面可能会被换到磁盘上，给其他页面腾地方。为了确
     保你正在使用的页面不会被换出去，必须先把它“锁定”。但这会增加一些额外的工作量。
在多处理器和多核系统中，所有这些开销都会变得更严重。表面上看，页表映射看起来非常高效——它
似乎可以在固定时间内完成，跟数据包大小没关系。但实际上，这个“固定时间”里的“固定成本”非常
高。 针对这个问题，Druschel和Peterson提出了一个叫做“fbufs”（快速缓冲区）的操作系统功能，成
功解决了大部分甚至全部的页面重映射开销问题。
fbufs的核心思想：
如果应用程序需要通过网络向操作系统内核发送大量数据包，那么某些缓冲区可能会被反复使用。操作
系统可以提前计算好（预计算）这些缓冲区的所有页面映射信息，等到实际传输数据时就不需要再花时
间映射页面了。或者，可以在第一次传输时建立这些映射（虽然前几个数据包会慢一些），但之后把这
些映射缓存起来，后续的数据包就不会有额外开销了。这样一来，在大多数情况下，页面重映射的开销
就被彻底消除了。
实现fbufs最简单的方法是使用共享内存。比如，把多个页面P1到Pk同时映射到内核和所有发送数据的应
用程序A1到Ak的页表中。但这其实是个糟糕的想法，因为这样可能会导致应用程序A1不小心读取到A2发
的数据包，破坏安全性和隔离性。
更好的办法是为每个应用程序到内核的传输单独保留（或懒惰地建立）一组共享页面。例如，可以为
FTP、HTTP等不同的服务分别设置一组缓冲区（页面）。更广泛地说，有些操作系统除了内核和应用程
序之外，还会定义多个安全子系统。因此，fbuf的设计者提出了一种叫做“路径”的概念，用来描述数据
在不同安全域之间的流动。比如，对于前面提到的简单场景，路径可以是从内核到应用程序，或者从应
用程序到内核（比如FTP到内核，或者内核到HTTP）。稍后我们会解释为什么路径是单向的—也就是
说，为什么每个应用程序都需要两条路径（一个方向一条）。



                                                        7
提前在路径中每个域的页表里映射好页面，或者等到需要时再创建缓冲区页面，这样可以避免在初始化
完成后实时处理路径时重新映射页面的额外开销。
上面的图展示了一个稍微复杂点的例子：以太网功能是内核级驱动程序实现的，TCP/IP协议栈是在用户
级的安全域中实现的，而Web应用程序则在应用层运行。每个安全域都有自己独立的一套页表。接收数
据的路径依次经过以太网、TCP/IP、Web，还有一条路径是以太网、OSI、FTP。
通过提前映射或按需懒加载的方式把缓冲区页面加到路径中每个域的页表里，可以在数据包第一次到达
并发送后，避免后续实时路径中重新映射页表的麻烦。
说到这里，你可能会问，为什么路径只能单向走？原因是，路径上的第一个进程被认为是写入者，后面
的进程都只负责读取。这可以通过设置页表条目中的权限位来实现：第一个进程的页表条目允许写操
作，而其他进程的页表条目只允许读操作。显然，这种方式是有方向性的，在两个方向上不对称，所以
路径必须是单向的。不过，这种设计确实提供了一定程度的保护。




                                               8
                      写优化
在上图中，我们用两个域的路径来展示这个过程。需要注意的是，在路径上的第一个应用程序里，页面8
被“预映射”到了第二个应用程序的页面10。这其实挺麻烦的，因为当第二个应用程序读取页面8的信息
时，它得想办法知道这个页面其实对应的是它自己这边的虚拟页面10（VP 10）。为了避免这种复杂性，
设计者遵循了一个简单的原则—避免不必要的复杂性，并规定fbuf在路径上的所有应用程序中都映射到
相同的虚拟页面（VP）。要实现这一点，可以在每个进程的虚拟内存（VM）中预留一些初始页面，专门
用来作为fbuf页面。
到目前为止，你可能觉得问题已经解决了，但实际上还有一些棘手的地方需要处理。为了保证安全性，
我们只允许一个程序负责写入数据，而其他程序只能读取数据。但这样一来，这些页面就变成“只读”的
了—只有写入者才能修改它们。那么问题来了：当数据沿着调用栈向下传递时，怎么给数据包添加头部
呢？解决方案如下图所示。这里的数据包实际上是一个组合结构，由多个部分组成，包括指向不同fbuf
的指针。通过向这个组合结构中添加普通缓冲区或者新的fbuf，就可以轻松地为数据包添加头部了。




                                                9
Using aggregate objects to allow adding layers to add headers while allowing only a single writer.
              使用聚合对象来允许添加层以添加头部信息，同时只允许单个写入器。

截至目前，fbuf方案已经采用了底层虚拟内存（VM）映射的思想，并通过分摊映射成本（尤其是在大规
模数据包传输场景中）来进一步提升效率。在典型情况下，页表更新操作被完全消除，这一优化可以在
常规操作系统中实现。如何在引入高效机制的同时保留标准的拷贝语义？具体而言，如果应用程序在将
一个fbuf传递给内核后继续对其进行写操作，系统应如何处理？传统的操作系统（如UNIX）无法依赖写
时复制（Copy-On-Write, COW）机制来解决这一问题。
fbufs最终的设计选择是放弃对标准拷贝语义的支持，转而对API进行修改。在这种新模型下，应用程序
开发者需要特别注意，避免在将fbuf传递给内核后对其执行写操作，直到该fbuf被内核归还至空闲列
表。为了防止因错误或恶意代码导致意外写入，内核会在fbuf从用户态传递至内核态时短暂地禁用写权
限（即切换写允许位），并在fbuf归还时重新启用该权限。如果应用程序尝试在写权限被禁用的情况下
进行写操作，将触发异常并导致应用程序崩溃，但其他进程的运行不会受到影响。
由于切换写允许位会引入一定的开销，而这正是fbufs试图避免的性能损耗，因此fbuf设施还提供了一种
替代形式的fbuf，称为“易失性（volatile）”fbuf。值得注意的是，如果写入方是一个可信实体（如内核
本身），则无需强制实施写保护。即使内核中存在导致意外写入的错误，整个系统的稳定性仍可能受到
影响，因为这种错误通常会导致系统级崩溃。
这种对API的改动乍看之下显得颇为激进。这是否意味着需要重写大量现有的基于UNIX网络栈的应用程
序？鉴于此任务的不可行性，研究者提出了多种解决方案。首先，可以通过扩展现有API的方式添加新的
系统调用。例如，Thadani和Khalidi（1995年）在其Solaris扩展中引入了uf_write()调用，作为标准
write()调用的补充。那些希望提升性能的应用程序可以选择使用这些新增的系统调用来重写部分逻辑，
从而利用fbuf机制带来的优化。
其次，这些扩展可以用于实现常见的I/O子系统（如UNIX的stdio库），这些子系统往往被多个应用程序
共享。通过链接到经过优化的库，应用程序本身无需任何修改即可从中受益。最终，实际的关键问题并
不在于API是否发生了变更，而在于应用程序开发者需要付出多大的努力来适应这些变更并从中获得性能
提升。
总结来说，fbuf方案通过牺牲标准拷贝语义和引入额外的开发约束，在特定场景下显著提升了性能。尽
                                                                                                10
管这种设计带来了API层面的重大变化，但通过合理的扩展和兼容性设计，其在实际应用中的推广障碍可
以得到有效缓解。
2.5 零拷贝在今天是否被使用？
随着我们转向100 Gbps，在进行了所有主要优化（例如，称为巨型帧的大MTU大小和一种称为分段卸载
的批处理形式，我们将在下一章中学习）之后，Linux网络栈的吞吐量在每个核心上饱和于约42 Gbps。
剩下的主要瓶颈仍然是内核与用户空间之间以及相反方向的数据拷贝。考虑在Linux中通过页表映射实现
零拷贝的提议的时机可能已经成熟，无论是在接收端（A reworked TCP zero-copy receive API，2018
年）还是在发送端（sendmsg，2022年）。
发送端的零拷贝提案由De Bruijn提出（sendmsg），使用了前几章描述的思想。再次假设应用程序必须
锁定并小心不要重用正在被内核使用且映射的缓冲区。应用程序可以通过内核放置在套接字相关错误队
列中的通知消息来判断何时可以安全重用缓冲区页面，序列号允许将通知与已完成的发送调用匹配。
接收端的提议更为复杂，因为内核必须将接收到的数据包放入通用缓冲区，确定其目标应用程序，然后
将缓冲区映射到应用程序的用户空间。正如我们之前在本文中看到的，这需要网络适配器或NIC将头部与
数据分离，并将数据对齐到页面。补丁（A reworked TCP zero-copy receive API，2018年）将MTU设
置为61,512——这允许十五个4096字节的页面数据，加上IPv6的40字节和TCP的32字节。这个补丁的有
趣之处在于，它不是使用内存映射（mmap()）来处理文件，而是实现了针对TCP套接字的mmap()。在
受控环境中的基准测试显示，每兆字节处理时间的潜在减少为3倍（A reworked TCP zero-copy
receive API，2018年）。
mmap是Linux提供的一种内存映射文件的机制，它实现了将内核中读缓冲区地址与用户空间缓冲区地址
进行映射，从而实现内核缓冲区与用户缓冲区的共享。这样就减少了一次用户态和内核态的CPU拷贝，
但是在内核空间内仍然有一次CPU拷贝。
请注意，这些最近的提议都没有使用本文中进一步优化的建议，例如fbufs来分摊页映射的成本，使其在
一系列调用中有效。此外，它们仅被考虑，并未在实际部署中使用。然而，零拷贝机制今天已广泛部
署，作为两种其他更专业技术的部分。第一种是通过DPDK（Data Plane Development Kit）（2018
年）方法进行的内核旁路，我们将在下一章中学习。第二种是涉及RDMA的系统。
3 使用远程DMA避免拷贝
虽然fbufs提供了一种减少应用程序和内核之间重复拷贝的合理方法，但其实还有更直接的办法，这种办
法不仅能减少拷贝，还能大幅降低控制上的开销。举个例子，如果你要通过以太网把一个1MB的文件从
一台工作站传到另一台，这个文件通常会被切成一堆小块，每块1460字节。然后，CPU得挨个处理这些
小块，完成TCP协议相关的工作，并且把每个数据包（可能通过类似fbufs的零拷贝接口）搬到应用程序
的内存里。
换个角度想想，CPU是怎么协调磁盘和内存之间的DMA（直接内存访问）操作的。CPU只需要告诉磁
                                                                11
盘：“嘿，把数据写到这块地址范围就行。”然后它就可以去干别的事了，等磁盘把一兆字节的数据都传
完，才会通知CPU：“主人，活儿干完了。” 请注意，这跟刚才网络传输的情况不一样，在那里CPU得对
每个小块数据亲力亲为。
这个对比启发了我们：为什么不能在网络上传输数据时也用类似DMA的方式呢？这就是所谓的RDMA
（远程直接内存访问）。实际上，这个想法最早是由一群计算机架构师在VAX Clusters系统中提出的，
时间可以追溯到1986年。有趣的是，据说当时有个搞VAX Clusters的人跑去DEC的网络团队学习网络知
识，结果被嘲笑了，还被扔了一本大学教材。但半年后，他带着RDMA的设计回来了。
RDMA的核心思想是：数据应该可以直接在两台计算机的内存之间通过网络传播，而不需要两边的CPU
对每个数据包都指手画脚。换句话说，两个网络适配器应该自己搞定一切：一个负责从内存里读数据，
另一个负责把数据写到目标内存里，整个过程就像DMA一样高效。不过，要实现这个目标，有两个关键
问题需要解决：
 1. 接收端的适配器怎么知道数据该放到哪？如果它还得回头问主机，那就违背了初衷。

 2. 安全性怎么保证？毕竟，恶意数据包可能会试图覆盖接收方的重要内存区域，这可不是小事。
    接下来，我们会先回顾一下这个早期的想法，然后再看看它在现代技术中的实现，比如光纤通道、
    Infiniband 方案。
3.1 在集群中避免拷贝
三十年前，数字设备公司（DEC）推出了一款非常成功的商业产品，叫VAX Clusters。它是一个可以扩
展计算能力的平台，专门为数据库应用等需求设计。这个系统的核心是一个叫计算机互连（CI）的高速
网络，速度达到140-Mbit，使用类似于以太网的协议。用户可以把多台VAX计算机和一些带网络功能的
磁盘连接到这个系统上。由于需要在远程磁盘和VAX内存之间传输大量数据，高效的数据拷贝变得非常
重要，也正是在这个背景下，RDMA技术应运而生。
RDMA的目标是让数据包一到达目标设备时，就能直接进入它们最终要去的地方。听起来简单，但实际
操作起来并不容易。在传统的网络中，数据包到达后，处理器至少得先检查一下这些数据包，然后决定
它们该去哪儿。即使CPU看了数据包的头部信息，也只能根据目标应用程序来决定把数据放到哪个接收
缓冲区队列里。
举个例子：假设某个接收应用程序已经为应用程序1分配了页面1、2和3，供适配器使用。如果第一个数
据包到达并被送到页面1，而第三个数据包先于第二个数据包到达，并且被错误地放到了页面2而不是页
面3，问题就来了。因为页面1、2和3是用来存储接收文件的，CPU可以在传输结束后重新整理这些页
面，但对于大文件来说，这种整理会非常麻烦。而且，即使是在理论上顺序传输的链路上，由于数据包
可能丢失或乱序，这种情况还是可能发生。




                                                     12
                       DMA across the network
VAX Clusters的做法是这样的：让目标应用程序提前锁定一组物理页面（比如图中的页面11和16），这
些页面就是文件传输的目标内存区域。不过，逻辑上这些物理页面会被显示成一组连续的逻辑页面（比
如图8中的页面1和2），统称为缓冲区B。这个缓冲区的名字B会被传给发送方应用程序。
接下来，发送方会在每个数据包的协议头中附带缓冲区名称和偏移量信息（比如P10）。这样一来，假如
第三个数据包没有按顺序到达，但它包含了缓冲区B和页面3的信息，那么这个数据包可以直接存到缓冲
区B的页面3中，即便它比第二个数据包先到。这样，等所有数据包都到达后，就不需要再花时间重新整
理页面了。
为了实现“不让处理器每次数据包到达时都分心”的目标，还需要满足几个额外的要求。首先，适配器必
须能自己处理传输协议（包括重复检查等工作），就像TCP那样。其次，适配器得能区分数据包的头部
和数据部分，只有数据部分才能被拷贝到目标缓冲区。
最后，允许任何携带缓冲区ID的数据包直接写入内存其实挺危险的，可能成为一个安全隐患。为了降低
风险，缓冲区ID里会包含一个难以猜测的随机字符串。更重要的是，VAX Clusters只在集群内部的可信
主机之间使用。如果要把这种方法推广到互联网上的数据传输，难度会大得多。
3.2 RDMA的现代实现
VAX Clusters引入了一个早期的存储区域网络（SAN）。存储区域网络（SAN）是连接多台计算机到共
享存储的后端网络，通常通过网络附加磁盘实现。今天有几个VAX Clusters的继任者提供SAN技术。
光纤通道(Fiber Channel)
1988年，美国国家标准协会（ANSI）下的一个工作组开始制定一种叫“光纤通道”（Fibre Channel）的
新标准。这个技术的目标之一是把当时常用的SCSI接口（小型计算机系统接口，用来连接工作站和本地
磁盘）扩展到更远的距离。所以，在很多光纤通道的使用场景中，SCSI协议仍然被用在光纤通道上。
相比VAX Clusters这种早期的网络技术，光纤通道在网络底层设计上更先进。它用了些现代技术，比如
通过交换机连接的点对点光纤链路。这让它的速度能飙到64 Gbps，还能覆盖比VAX Clusters大得多的
距离。交换机甚至可以远程连接，这让一些交易公司能在远处备份所有交易数据。不过，用了交换机就
                                                       13
得注意流控问题，特别是要小心避免丢包，这样才能保持网络看起来像一个无损的计算机总线。
在安全性方面，光纤通道比VAX Clusters稍微强一点。在VAX Clusters里，只要设备名字对上了，就能
随便访问其他设备的存储。而光纤通道可以把网络分成多个“区域”，每个区域的设备只能看到自己区域
内的存储，看不到其他区域的东西。一些新产品还引入了认证机制，进一步加强了安全。
尽管底层技术有这些不同，但核心思想其实差不多。通过命名缓冲区进行远程直接内存访问（RDMA）
仍然是关键技术。虽然光纤通道在和基于以太网的RDMA技术（比如RoCE）竞争时面临挑战，但它依然
是存储区域网络（SAN）的核心技术，并且短期内不会被淘汰。
光纤通道之所以在SAN中一直受欢迎，可能和固态硬盘（SSD）以及非易失性存储器（NVM，比如PCI
Express上的NVMe）的普及有关。光纤通道的延迟特别低，这对高性能存储非常重要。虽然以太网方案
在成本上有一定优势，但光纤通道的性能可能是它在SAN中仍然占据主导地位的原因。
Infiniband
InfiniBand的设计初衷源于对当时工作站和PC内部I/O总线（如PCI总线）性能瓶颈的观察。尽管PCI总线
的最大带宽为533 MB/s，但随着现代高速外设（例如千兆以太网接口卡）的普及，其性能已无法满足需
求。尽管市场上出现了一些临时性解决方案（如PCI-X总线），但内部计算机互连技术亟需一种类似于外
部网络从10-Mbit以太网向千兆以太网演进的扩展能力。
此外，研究发现，计算机内部通常存在三种独立的网络技术：用于网络连接的接口（如以太网）、用于
存储设备的接口（如通过光纤通道实现的SCSI）以及用于内部设备通信的PCI总线。基于奥卡姆剃刀原
则，简化系统架构的需求促使业界考虑用一种统一的网络技术替代这三种异构技术。在此背景下，康
柏、戴尔、惠普、IBM和Sun等公司联合成立了InfiniBand贸易协会，旨在推动这一目标的实现。
InfiniBand规范借鉴了许多光纤通道底层网络技术的核心思想，并采用基于交换机和点对点链路的架构设
计。与此同时，InfiniBand引入了多项创新特性：它以下一代互联网协议中的128位IP地址为基础进行寻
址；支持将单个物理链路虚拟化为多个独立的逻辑链路（称为通道）；具备服务质量（QoS）和多播功
能；并采用了远程直接内存访问（RDMA）技术，显著减少了数据传输过程中的内存拷贝开销。
尽管InfiniBand未能像以太网解决方案那样迅速普及，也未在存储区域网络（SAN）领域取代光纤通道的
地位，但它在高性能计算（HPC）领域依然占据重要地位。特别是在超级计算环境中，InfiniBand被广泛
应用于构建高性能互连网络（如SHARP技术，2019年）。此外，InfiniBand还被用于加速大规模机器学
习工作负载，通过网络计算抽象技术（如SHARP）提升分布式计算效率。例如，在橡树岭国家实验室和
洛斯阿拉莫斯国家实验室部署的全球最大的两台超级计算机中，每台均使用了近10,000个InfiniBand节点
（SHARP，2019年）。
RoCE 融合以太网的RDMA：
Infiniband中的RDMA（远程直接内存访问）运行在一种使用基于信用的流控制机制的网络上，这种机制
的主要目的是确保网络无损。通过这种方式，Infiniband能够在高性能计算和数据中心环境中提供极高的
吞吐量和极低的延迟。然而，由于Infiniband的设计初衷是运行在几乎无丢包的网络环境中，因此其
RDMA传输协议（通常在网卡NIC上实现）并未被优化以高效地处理数据包丢失的情况。当接收方接收到
乱序的数据包时，它会直接丢弃该数据包，并向发送方发送一个NACK（否定确认）。这种行为会导致发
                                                        14
送方重新传输整个滑动窗口中的所有数据包，从而显著降低性能，特别是在出现丢包的情况下。
相比之下，RoCE（基于融合以太网的RDMA）的核心思想是将RDMA功能移植到标准以太网上，而不是
依赖于Infiniband专用硬件。RoCE选择在UDP之上运行，而非TCP，这使得它能够继承UDP的轻量级特
性，同时避免了TCP的复杂连接管理和拥塞控制开销。然而，为了防止因网络拥塞而导致的丢包（这是
以太网中最常见的丢包原因），RoCE引入了一种称为PFC（优先流控制）的粗略形式的流控制机制。
PFC的作用是在交换机层面进行流量管理，当下游交换机感知到拥塞时，它会向上游交换机发送信号，
要求其暂停特定优先级的数据流，直到下游交换机的缓冲区压力得到缓解。通过这种方式，PFC可以在
数据包实际被丢弃之前采取预防措施，从而减少丢包的可能性。
需要指出的是，PFC这种流控制机制与Infiniband中基于信用的流控制相比，虽然更加简单易实现，但它
也存在一定的局限性。例如，PFC的“暂停”信号可能会导致整个网络中的流量暂停传播，从而引发所谓
的“暂停风暴”问题。此外，PFC只能针对特定优先级的流量进行控制，而无法像Infiniband那样对每个数
据流进行精细化的信用管理。尽管如此，PFC作为一种折衷方案，能够在现有的以太网基础设施上实现
接近无损的网络环境，这对于RoCE的广泛应用至关重要。总的来说，这两种技术分别代表了两种不同的
设计理念：Infiniband追求极致的性能和精确的流控制，而RoCE则更注重兼容性和易部署性。
\subsection{案例研究：数据中心拥塞控制的演进}

传统TCP的拥塞控制基于丢包信号——当检测到丢包时降低发送速率。这种"丢包即降速"的粗放策略在现代数据中心面临严峻挑战：AI训练和分布式存储等应用对长尾延迟极度敏感，一个数据包的延迟可能导致成千上万个GPU的空转。过去十年，SIGCOMM上的工业界研究展示了拥塞控制从"反应式"向"主动式"、从"粗放"向"精确"的演进。

\textbf{微软Azure与DCQCN}：微软在SIGCOMM 2016发表的RoCEv2大规模部署研究是网络领域最具影响力的经验分享之一。RDMA能够绕过CPU直接操作远程内存，但其对"无损网络"的苛刻要求使其在超大规模以太网部署中充满风险。

微软通过引入基于DSCP的优先级流量控制（PFC），成功解决了原有VLAN-based PFC在三层路由环境下的局限性。为了控制PFC可能引发的死锁和"暂停风暴"，微软开发了\textbf{数据中心量化拥塞通知（DCQCN）}算法。DCQCN利用交换机的显式拥塞通知（ECN）信号，结合端点的速率限制器，实现了平滑的流量调节。

微软的实践证明，通过精细的参数调优和硬件协同，原本属于实验室环境的RDMA可以在包含上万台服务器的公共云中稳定运行，将99.9百分位的网络延迟从毫秒级降低到了百微秒级。

\textbf{谷歌Swift协议}：谷歌于2020年提出的Swift协议针对云环境下的多租户和高动态性进行了优化。Swift将RTT（往返时间）作为主要的拥塞信号，通过对\textbf{亚微秒级延迟变化}的精确捕捉，在丢包发生前就主动调整发送速率。

这种"延迟敏感"的特性对于分布式存储和RPC调用至关重要，因为它能有效抑制长尾延迟（Tail Latency），从而保证上层应用的SLA。Swift的设计哲学代表了拥塞控制的新方向：从"丢包后补救"到"延迟前预防"。

\textbf{AWS SRD协议}：亚马逊在2020年推出的可扩展可靠数据报（SRD）协议做出了一个大胆的选择：提供\textbf{可靠但无序的交付}。SRD发现，在现代多路径数据中心网络中，TCP的按序交付约束和对丢包的敏感性已成为性能瓶颈。

SRD的核心机制包括：\textbf{多路径并发发送}——利用ECMP路径多样性将数据包随机喷淋到成百上千条路径；\textbf{硬件级微秒重传}——在定制的Nitro网络卡中实现，能够在微秒级感知链路波动并进行重传；\textbf{乱序容忍}——主要服务于EBS存储和HPC/AI负载，这些应用可以通过应用层缓存处理乱序，从而换取极低的尾部延迟。

通过将SRD协议封装在EFA接口中，AWS成功将原本只能在专用超级计算机上运行的MPI和NCCL负载迁移到了公共云中，构建了包含近50万颗Trainium芯片的超大规模AI算力集群。

\textbf{趋势总结}：从DCQCN到Swift再到SRD，拥塞控制的演进呈现三个共同趋势：\textbf{更细粒度的信号}（从丢包到ECN再到亚微秒RTT）、\textbf{更主动的干预}（从反应式降速到预防式调度）、\textbf{更深度的软硬协同}（算法与DPU/智能网卡紧密集成了）。这些创新使得数据中心网络能够支撑AI时代对带宽和延迟的极致需求。

4 扩展到文件系统




回想一下上图，为了处理对文件X的请求，服务器可能需要从磁盘读取X（Copy 1）到内核缓冲区（代表
文件缓存），然后将文件缓存拷贝到应用程序缓冲区（Copy 2）。如果文件已在缓存中，Copy 1在服务
器有足够内存的情况下通常会被忽略，这对于流行文件来说是合理的假设。主要目标是移除Copy 2。请
注意，在Web服务器中不必要地将拷贝次数翻倍不仅会减半有效总线带宽，还可能减半服务器缓存的大
小。这反过来会由于更高的未命中率而降低服务器性能，这意味着更大比例的文档以磁盘速度而非总线
速度提供服务。

                                                     15
本节调查三种用于移除冗余文件系统拷贝（Copy 2）的技术。首先，描述一种称为共享内存映射的技
术，可以减少Copy 2，但与网络子系统集成不佳。其次，描述IO-Lite，本质上是将fbufs扩展到包括文
件系统。最后，描述一种称为I/O拼接的技术，许多商业Web服务器都在使用。
4.1 共享内存
现代UNIX系统（Stevens，1998年）提供了一个方便的功能叫 mmap()，它可以让程序（比如服务器）
把文件直接“映射”到自己的内存空间里。其他操作系统也有类似的功能。简单来说，当一个文件被映射
到程序的内存中时，就好像这个程序在内存里有了文件的一个副本。这听起来好像有点多余，因为文件
系统本身已经在内存里缓存了文件的内容。不过，由于操作系统的虚拟内存机制很聪明，实际上并没有
真正复制文件内容，而是让多个程序和文件系统共享同一份物理内存中的文件数据。
有一个叫Flash Web服务器的软件（Pai等人，1999年），它利用 mmap()把经常用的文件映射到内存
中，从而避免了图1中提到的两次数据拷贝（Copy 1和Copy 2）。但由于内存资源有限，Flash不能把所
有文件都一直放在内存里，而是只能缓存文件的一部分。它会用一种叫LRU（最近最少使用）的策略，
把长时间没用过的文件从内存中移除。
需要注意的是，文件系统本身其实已经有类似的缓存功能了（而且它运行在内核里，对内存资源的管理
更精准）。但Flash之所以还要自己搞一套缓存，是为了避免那两次不必要的拷贝。虽然Flash通过 mmap
()绕过了文件系统的拷贝，但它仍然受限于UNIX系统的API，因此在网络传输时还是避免不了第三次拷
贝（Copy 3）。就在我们以为解决了Copy 2的问题时，讨厌的Copy 3又冒出来了！
有一种方法可以同时避免Copy 3，那就是结合使用 TCOW和 mmap()（Brustoloni和Steenkiste，
1996年）。但这也有它的缺点，而且这种方法并不通用，无法解决通过UNIX管道与CGI进程交互时的数
据拷贝问题。
4.2 IO-Lite：缓冲区的统一视图
虽然把仿真拷贝和 mmap()结合起来可以去掉所有多余的数据拷贝，但这种方法还是有一些不足的地
方。首先，它没法避免CGI应用程序和Web服务器之间的数据拷贝。这些CGI程序通常是以独立进程运行
的，它们通过UNIX管道把动态内容传给服务器进程。而管道这种进程间通信方式，本质上还是需要在两
个地址空间之间拷贝数据，这就带来了额外的开销。
其次，我们发现现有的方案都没解决一个很耗时的操作——TCP校验和计算。当同一个文件被缓存命中
时，除了第一个包含HTTP头的响应包外，后面返回文件内容的数据包对所有请求来说都是完全一样的。
那为什么不能把这些TCP校验和也缓存起来呢？这确实是个好主意，但传统缓冲方案效率太低，因为要
建立一种能根据数据包内容映射到校验和的缓存机制。
这一节介绍的IO-Lite缓冲方案（Pai等，1999b）是基于fbuf思想的一个扩展，专门应用到了文件系统领
域。它的特点是可以去掉上图中所有的冗余拷贝，甚至还能消除CGI进程和服务器之间的重复拷贝。它通
过一种特殊的缓冲区编号机制，让子系统（比如TCP）能够快速识别哪些数据包是重传的，从而进一步
优化性能。

                                                               16
IO-Lite可以看作是fbuf的“升级版”，但由于整合了文件系统，复杂度明显提高了。需要注意的是，fbuf
和 mmap()没法一起用。这是因为 mmap()允许应用程序自己选择缓冲区的地址和格式，而fbuf的地
址和格式是由内核决定的。所以，如果应用程序用虚拟地址空间映射了一个文件，想直接通过fbuf发送
数据就不行了，因为两者的地址空间不一致，必须得做一次物理拷贝。为了解决这个问题，IO-Lite把
fbuf的概念扩展了一下，覆盖了文件系统的场景，这样就不需要依赖 mmap()了。另外，IO-Lite是在
通用操作系统（比如UNIX）上实现的，这点和fbuf不一样。不过除了这两点区别，IO-Lite继承了fbuf的
所有核心思想：通过不可变的缓冲区（IO-Lite里叫slices）实现只读共享、使用复合缓冲区（叫buffer
aggregates），以及针对路径延迟创建缓冲区缓存（叫I/O stream）。
尽管核心思想差不多，但IO-Lite还得面对一些新挑战，主要来自和文件系统的整合：
  1. 它需要处理多个应用、TCP代码和文件服务器之间复杂的共享缓冲区模式。

  2. IO-Lite的页面可能同时扮演两种角色：一种是虚拟内存页面（由磁盘分页支持的文件备份），另一
     种是文件页面（由实际磁盘文件支持）。
  3. 它得找到一种在UNIX系统中优雅集成的方式，同时尽量避免大规模改动现有系统。

目前，AWS和Azure等云计算厂商在高性能存储服务（如AWS Nitro SSD）中部分采用了IO-Lite
的零拷贝思想。阿里云在2023年发布的“神龙架构”升级中，通过定制网卡支持类似IO-Lite的共享
内存传输，用于其分布式数据库PolarDB的日志同步。




通过有效地传递指向单一IO-Lite缓冲区的指针（通过虚拟内存映射），IO-Lite消除了所有冗余拷贝。假
设文件、TCP校验和以及HTTP响应都已缓存，Web服务器只需要将这些缓存值一次性拷贝到网络接口即
可。上图展示响应GET请求的新流程：文件从磁盘读入系统缓存时以IO-Lite缓冲区形式存储；应用读取
文件时不进行物理拷贝，而是创建指向该缓冲区的buffer aggregate；当应用发送文件至TCP时，网络系
统获取同一IO-Lite页面的指针。IO-Lite通过引用计数确保缓冲区安全重用。

                                                      17
该方案还包含两项优化：1）应用缓存常见文件的HTTP响应，通常只需最小修改即可复用；2）IO-Lite
为每个缓冲区分配唯一编号，TCP模块维护按编号索引的校验和缓存，使得重复传输文件时只需首次计
算校验和。这些改进彻底消除了冗余操作。
IO-Lite还可实现零拷贝管道：当CGI进程与服务器进程采用该IPC机制时，在保持进程隔离安全性的前提
下消除所有拷贝。此外允许应用定制缓存策略，支持Web服务器基于大小和访问频率的智能缓存。值得
注意的是，IO-Lite无需完全移除内核或深度绑定应用与内核即可实现这些优化。
4.3 通过I/O拼接避免文件系统拷贝




      图片加载失败
在I/O拼接中，通过一个“拼接”磁盘I/O流和网络I/O流的系统调用，消除了与用户空间缓冲区之间拷贝带
来的所有间接操作。一如既往，如果文件在缓存中，则可以省略拷贝操作1。

I/O拼接是一种常用的技术，目的是减少文件系统和网络子系统之间不必要的数据拷贝。简单来说，它的
核心思想就是把多个小的读写操作合并成一个大的操作，从而减少系统调用的次数和数据拷贝的开销。
具体怎么实现呢？I/O拼接通过维护一个大的缓冲区池来完成任务。比如，当应用程序需要读取文件时，
Web服务器会把多个小文件块先合并到一个大的缓冲区里，然后一次性把这个缓冲区的数据传给网络子
系统。同样，当应用程序要发送数据时，服务器也会先把多个小数据块合并到一个大缓冲区，再一次性
通过网络发送出去。
这种方式的好处很明显：它减少了系统调用的次数，因为每次系统调用都会消耗一定的资源。同时，通
过把小数据块合并成大数据块，也减少了数据拷贝的次数——数据只需要从文件系统拷贝到网络缓冲区
一次，而不是多次。
I/O拼接的实现通常依赖于操作系统提供的一些底层机制，比如 mmap()和共享内存。有了这些机制，
Web服务器可以直接访问文件系统和网络缓冲区，省去了传统意义上的数据拷贝过程。
                                                  18
尽管I/O拼接很有效，但它也有一些缺点。首先，它需要应用程序和操作系统紧密配合，这可能限制了它
的通用性。其次，合并小的操作可能会增加延迟，因为应用程序得等足够多的数据准备好才能继续下一
步。最后，管理这些缓冲区也需要精心设计，否则可能会导致缓冲区耗尽或者性能瓶颈。
不过，尽管有这些问题，I/O拼接仍然是一项非常有效的优化技术，尤其是在高负载的Web服务器环境
中。它证明了通过优化I/O操作可以显著提升系统性能，并为未来的研究和开发提供了重要参考。
远程直接内存访问（RDMA）与文件系统的结合
随着RDMA技术越来越成熟，研究人员开始琢磨怎么把它和文件系统结合起来，进一步提升I/O性能。简
单来说，RDMA能让数据直接在内存之间传输，完全不用CPU插手。这为解决文件系统和网络子系统之
间的重复拷贝问题提供了新思路。
一种常见的做法是把RDMA和共享内存机制结合起来。具体来说，文件系统把数据存放在共享内存区
域，而RDMA可以直接访问这些区域，这样就省去了传统方式中那些繁琐的数据拷贝步骤。这种方法特
别适合高性能计算（HPC）环境，因为在这些场景下，多个节点需要频繁地共享和交换数据。
还有一种方法是利用RDMA的“零拷贝”特性来优化文件系统的缓存机制。在这种方式下，文件系统的缓
存直接放在RDMA能访问的内存区域，远程节点可以直接读取这些缓存，既减少了延迟，又降低了数据
传输的开销。这种方式在高负载的Web服务器环境中也有很大的应用潜力。
不过，尽管RDMA和文件系统的结合看起来很有前景，实际操作时还是有不少挑战。首先，RDMA需要专
门的硬件支持，这可能会增加系统成本。其次，安全性是个大问题，必须设计得当，防止未经授权的访
问和数据泄露。最后，RDMA的管理需要和文件系统的管理紧密结合，才能保证数据的一致性和完整
性。
即便如此，RDMA与文件系统的结合仍然是一个值得研究的方向，特别是在那些对带宽要求高、延迟要
求低的应用场景中。它展示了通过整合不同技术来优化I/O性能的可能性，也为未来的研究和开发提供了
重要参考。
5 总结
与线性规划中的椭球算法和快速混合马尔可夫链理论等理论技术的精妙和复杂相比，诸如避免拷贝等系
统技术似乎平淡无奇。然而，可以认为系统的复杂性不在于深度（即每个组件自身的复杂性），而在于
广度（即组件之间的复杂关系）。或许，优化Web服务器内存带宽所需的对HTTP、文件系统、网络代
码、指令缓存实现等方面的广泛理解，为这一观点提供了一些证据。




                                              19
